\section{Introduction}
\label{sec:intro}

Evolutionary search improves a solution by repeating two steps. A \emph{variation} operator makes
new candidates from old ones, and a \emph{selection} step keeps the better ones
\citep{back1997handbook}. For decades the variation operator was hand-designed: bit flips,
subtree swaps and crossover \citep{oneill2010openissues}. Large language models (LLMs) changed this. In
ELM \citep{lehman2022elm}, FunSearch \citep{romeraparedes2024funsearch}, EvoPrompting
\citep{chen2023evoprompting} and AlphaEvolve \citep{novikov2025alphaevolve}, an LLM writes the new
candidate. But a fixed framework still samples the parents, runs the evaluation and manages the
population. The LLM fills one slot in a pipeline:
\begin{equation}
  \Vary(\Pt) = \Generate\big(\Sample(\Pt)\big).
  \label{eq:intro-generate}
\end{equation}

The AVO paper \citep{chen2026avo} argues that this slot is too small. An LLM confined to
$\Generate$ ``produces a single output per invocation, with no ability to proactively consult
reference materials, test its changes, interpret feedback, or revise its approach before committing
a candidate'' \citep[\S1]{chen2026avo}. AVO gives the whole variation step to an autonomous coding
agent:
\begin{equation}
  \Vary(\Pt) = \Agent(\Pt, \K, \f),
  \label{eq:intro-agent}
\end{equation}
where $\Pt$ is the lineage of solutions and their scores, $\K$ is a domain knowledge base and $\f$
is the scoring function. The agent decides what to read, what to edit and when to test. NVIDIA ran
this operator for seven days on multi-head attention for Blackwell B200 GPUs. It committed 40
versions and reached up to \SI{1668}{TFLOPS}, which is up to 3.5\% faster than cuDNN and up to
10.5\% faster than FlashAttention-4 on the tested configurations.

The paper's result is about kernels. Its \emph{harness} is not. Equation~\eqref{eq:intro-agent}
needs only four things: a way to measure a candidate, a record of the versions that survived, a
place to look things up, and an agent. None of them is specific to GPUs. But the paper releases no
code, names no model and gives no supervisor thresholds. It leaves the acceptance rule for a
vector score as the phrase ``matches or improves''. A reader who wants to apply AVO to a different
problem has to rebuild the harness from the prose.

This article describes such a rebuild. \emph{avocode} is a command-line tool (\cmd{avo}) and a set
of agent skills that turn any git repository with a measurable score into an AVO loop. Its thesis
is that the harness generalizes past attention kernels to any repository with a strict, cheap
scorer. The design follows one rule throughout: the agent proposes, and the measurement decides.

\paragraph{Contributions.}
\begin{enumerate}[leftmargin=*,itemsep=2pt]
  \item \textbf{A concrete mapping of AVO to a repository} (\cref{sec:overview}). $\f$ becomes an
    executable with a one-line JSON contract, $\Pt$ becomes git commits with score trailers, and
    $\K$ becomes a hybrid-search index whose corpus includes the lineage itself.
  \item \textbf{A precise commit rule} (\cref{sec:lineage}). We define acceptance as weak Pareto
    dominance over the score vector with strict gain on at least one configuration and a
    symmetric relative noise floor. We state the invariants this rule gives and the ways it
    resists gaming.
  \item \textbf{An external, deterministic supervisor} (\cref{sec:supervise}). It detects two
    signals, stall and thrash, and emits a directive that \emph{cites} prior versions, recorded
    dead ends and documents in $\K$ that no version has used.
  \item \textbf{Agent-agnostic delivery} (\cref{sec:adapters}). The same five skill files drive
    pi, Claude Code and Codex. Every capability is reachable from a shell.
  \item \textbf{A bounded fan-out operator} (\cref{sec:fan}). \avo{fan} runs $N$ probes in
    disposable git worktrees on a small model. It is a first, local step toward the population
    search that the paper defers.
  \item \textbf{An empirical case study and an honest list of limits} (\cref{sec:eval,sec:limits}).
    Two runs on a fuzzy search benchmark reach a $7837\times$ speedup with no committed
    regression. We also show where the evidence is thin.
\end{enumerate}

\paragraph{Outline.} \Cref{sec:background} places AVO among related work. \Cref{sec:avo}
summarizes the AVO paper. \Cref{sec:overview,sec:score,sec:lineage,sec:orchestration} describe
avocode, from the system view down to the commit rule and the loop. \Cref{sec:eval} reports
results. \Cref{sec:discussion} compares avocode with the paper, and \cref{sec:limits} lists
limitations and future work.
